\documentclass{article}
\usepackage[T1]{fontenc}
\usepackage{lmodern}
\usepackage[a4paper, landscape, margin=14mm]{geometry}
\usepackage{xcolor}
\usepackage{tikz}

% Pick the month and year. The grid, weekdays and day count update automatically.
\newcommand{\calyear}{2026}
\newcommand{\calmonth}{10}

% Add events as \event{day}{text}. Use \\ for a line break.
\newcommand{\event}[2]{\expandafter\def\csname calevent#1\endcsname{#2}}
\event{2}{Team lunch}
\event{9}{Project\\deadline}
\event{15}{Dentist 10:00}
\event{21}{Book club}
\event{31}{Halloween party}

% Colours for the header, weekends and events.
\definecolor{calink}{HTML}{283044}
\definecolor{calaccent}{HTML}{EF8354}
\definecolor{calweekend}{HTML}{F1F3F6}

\renewcommand{\familydefault}{\sfdefault}
\pagestyle{empty}
\setlength{\parindent}{0pt}

\newcommand{\monthname}[1]{\ifcase#1\or January\or February\or March\or April\or May\or June\or July\or August\or September\or October\or November\or December\fi}

\begin{document}

% Weekday of the 1st (Monday = 0), days in the month and rows needed.
\pgfmathtruncatemacro{\yy}{\calmonth < 3 ? \calyear - 1 : \calyear}
\pgfmathtruncatemacro{\sunfirst}{mod(\yy + div(\yy,4) - div(\yy,100) + div(\yy,400) + {0,3,2,5,0,3,5,1,4,6,2,4}[\calmonth-1] + 1, 7)}
\pgfmathtruncatemacro{\offset}{mod(\sunfirst + 6, 7)}
\pgfmathtruncatemacro{\leap}{(mod(\calyear,4) == 0 && mod(\calyear,100) != 0) || mod(\calyear,400) == 0}
\pgfmathtruncatemacro{\ndays}{{31,28,31,30,31,30,31,31,30,31,30,31}[\calmonth-1] + (\calmonth == 2 ? \leap : 0)}
\pgfmathtruncatemacro{\nrows}{ceil((\offset + \ndays) / 7)}
\pgfmathsetlengthmacro{\cellh}{142mm / \nrows}

\begin{tikzpicture}[x=38.4mm, y=-\cellh]
  % Header with month and year.
  \node[anchor=south west, font=\bfseries\fontsize{44}{48}\selectfont, text=calink, inner sep=0pt]
    at (0,-0.1) {\monthname{\calmonth}};
  \node[anchor=south east, font=\fontsize{44}{48}\selectfont, text=calaccent, inner sep=0pt]
    at (7,-0.1) {\calyear};

  % Weekday names row.
  \foreach \name [count=\i from 0] in {Monday,Tuesday,Wednesday,Thursday,Friday,Saturday,Sunday} {
    \fill[calink] (\i,0) rectangle ++(1,-8mm);
    \node[text=white, font=\small\bfseries] at (\i+0.5,-4mm) {\name};
  }

  \begin{scope}[yshift=-10mm]
    % Weekend shading and grid.
    \fill[calweekend] (5,0) rectangle (7,\nrows);
    \foreach \c in {1,...,6} {\draw[calink!30] (\c,0) -- (\c,\nrows);}
    \foreach \r in {1,...,\nrows} {\draw[calink!30] (0,\r) -- (7,\r);}
    \draw[calink, line width=1pt] (0,0) rectangle (7,\nrows);

    % Day numbers and events.
    \foreach \d in {1,...,\ndays} {
      \pgfmathtruncatemacro{\pos}{\offset + \d - 1}
      \pgfmathtruncatemacro{\col}{mod(\pos, 7)}
      \pgfmathtruncatemacro{\row}{div(\pos, 7)}
      \node[anchor=north west, font=\large\bfseries, text=calink] at (\col+0.03,\row+0.04) {\d};
      \ifcsname calevent\d\endcsname
        \node[anchor=south west, fill=calaccent!20, draw=calaccent, rounded corners=2pt,
          font=\scriptsize, align=left, text width=30mm] at (\col+0.05,\row+0.95) {\csname calevent\d\endcsname};
      \fi
    }
  \end{scope}
\end{tikzpicture}

\vfill
{\small\color{calink!70} Notes:\ \hrulefill\par}

\end{document}
